% !TeX root = ../../Book5.tex
% 纲要标题：### 11.2 子代数、理想与商
% 纲要位置：工作记录/计划纲要.md，第 4427 行。

\section{子代数、理想与商}
\label{b5:sec:1102}


% 纲要范围：
%
% * 同态及同构定理
% * 中心与导子
% * 伴随表示
% * 区分两种理想，说明正规化子的最大性、导子的一阶来源及伴随像对中心信息的遗失
%

Lie 代数的子对象和商仍须与括号相容。理想条件保证一个括号在取商后不依赖代表；伴随表示则把这种相容性写成线性算子的交换关系。这些构造是后面分解 Lie 代数的基础。

\subsection{同态及同构定理}
\label{b5:subsec:1102:01}

\begin{definition}\label{b5:def:lie-subalgebra-ideal}
设 \(\mathfrak g\) 为域 \(k\) 上的 Lie 代数。子空间 \(\mathfrak h\subseteq\mathfrak g\) 若满足 \([\mathfrak h,\mathfrak h]\subseteq\mathfrak h\)，就称为子代数；子空间 \(\mathfrak a\subseteq\mathfrak g\) 若满足 \([\mathfrak g,\mathfrak a]\subseteq\mathfrak a\)，就称为理想。这里 \([U,W]\) 表示所有 \([u,w]\)、\(u\in U,w\in W\) 的线性生成空间。
\end{definition}
结合代数的双边理想满足比 Lie 理想更强的条件。若 \(A\) 为域 \(k\) 上的结合代数、\(I\) 为其双边理想，则对 \(a\in A\)、\(u\in I\)，两个乘积 \(au,ua\) 都在 \(I\) 中，因而 \([a,u]=au-ua\in I\)；所以 \(I\) 是交换子 Lie 代数 \(A^{-}\) 的理想。反向不成立：在任意域 \(k\) 上，\(\mathfrak{sl}_2(k)\) 是 \(\mathfrak{gl}_2(k)=M_2(k)^{-}\) 的 Lie 理想，因为 \(\operatorname{tr}[X,Y]=0\)。但
\[
E_{12}\in\mathfrak{sl}_2(k),\qquad
E_{21}E_{12}=E_{22}\notin\mathfrak{sl}_2(k),
\]
故它不是 \(M_2(k)\) 的双边理想。这里 \(\operatorname{tr}E_{22}=1\)，所以特征二时反例也成立。Lie 理想只要求乘积之差留在子空间，并不要求两种乘积分别留在其中。
\begin{definition}\label{b5:def:lie-homomorphism}
设 \(\mathfrak g,\mathfrak l\) 为同一域 \(k\) 上的 Lie 代数。一个 \(k\) 线性映射 \(f:\mathfrak g\to\mathfrak l\) 若满足
\(f\bigl([x,y]\bigr)=\bigl[f(x),f(y)\bigr]\) 对所有 \(x,y\in\mathfrak g\) 成立，就称为 Lie 代数同态；若它还是线性同构，就称为 Lie 代数同构。
\end{definition}
例如上三角矩阵在 \(\mathfrak{gl}_n(k)\) 中为子代数，严格上三角矩阵在上三角代数中为理想。后一空间通常不是整个 \(\mathfrak{gl}_n(k)\) 的理想：在二维中，
\([E_{21},E_{12}]=E_{22}-E_{11}\)
不再严格上三角。理想的环境不能从叙述中省去。
\begin{proposition}\label{b5:prop:lie-isomorphism}
设 \(\mathfrak g,\mathfrak l\) 为同一域 \(k\) 上的 Lie 代数，\(f:\mathfrak g\to\mathfrak l\) 为 Lie 代数同态。则 \(\ker f\) 为 \(\mathfrak g\) 的理想，\(\operatorname{im}f\) 为 \(\mathfrak l\) 的子代数，且
\[
\mathfrak g/\ker f\cong\operatorname{im}f.
\]
更一般地，对 \(\mathfrak g\) 的理想 \(\mathfrak a\)，商空间上的括号
\([x+\mathfrak a,y+\mathfrak a]=[x,y]+\mathfrak a\)
良定义；包含 \(\mathfrak a\) 的理想与商 Lie 代数的理想一一对应。
\end{proposition}
\begin{proof}
若 \(x\in\ker f\)，则 \(f\bigl([y,x]\bigr)=\bigl[f(y),0\bigr]=0\)，证明核为理想；像对括号封闭直接由同态等式得到。若更换两个商代表，括号的差为
\([a,y]+[x,b]+[a,b]\)，其中 \(a,b\in\mathfrak a\)，各项都在理想中，故良定义。双线性、交替性与 Jacobi 从原代数下降。映射 \(x+\ker f\mapsto f(x)\) 线性双射并保持括号，得到同构。最后，理想沿商映射取像或逆像时分别保持括号相容条件，两种操作互逆。
\end{proof}
对上述 \(\mathfrak{gl}_2(k)\)，迹映射到一维交换 Lie 代数 \(k\) 是 Lie 同态，因为交换子的迹为零。它的核为 \(\mathfrak{sl}_2(k)\)，且 \(E_{11}\) 的迹为 \(1\)，所以像为全部 \(k\)。将这组核与像代入\cref{b5:prop:lie-isomorphism}，得到 Lie 代数同构
\[
\mathfrak{gl}_2(k)/\mathfrak{sl}_2(k)\cong k.
\]
这个商的括号为零；原矩阵乘法却不能在同一个商空间上诱导结合代数乘法，因为所除的空间不是双边理想。

\ProseParagraphBreak

同样的计算还给出第二、第三同构定理。例如若 \(\mathfrak h\) 为子代数、\(\mathfrak a\) 为理想，则
\(\mathfrak h/(\mathfrak h\cap\mathfrak a)\cong(\mathfrak h+\mathfrak a)/\mathfrak a\)；
若 \(\mathfrak a\subseteq\mathfrak b\) 为理想，则
\((\mathfrak g/\mathfrak a)/(\mathfrak b/\mathfrak a)\cong\mathfrak g/\mathfrak b\)。
这些映射分别由包含和复合商映射给出，核与像已由前一命题确定。

\subsection{中心与导子}
\label{b5:subsec:1102:02}

\begin{definition}\label{b5:def:lie-center-derivations}
设 \(\mathfrak g\) 为域 \(k\) 上的 Lie 代数。其中心为
\[
Z(\mathfrak g)=\bigl\{\,z\in\mathfrak g\mid[z,x]=0\;\text{对所有}\;x\in\mathfrak g\,\bigr\}.
\]
\(k\) 线性映射 \(D:\mathfrak g\to\mathfrak g\) 若对所有 \(x,y\in\mathfrak g\) 满足
\[
D\bigl([x,y]\bigr)=\bigl[D(x),y\bigr]+\bigl[x,D(y)\bigr],
\]
就称为 \(\mathfrak g\) 的\term[lie-daishu-daozi]{导子}[derivation of a Lie algebra]。全部导子的空间记为 \(\operatorname{Der}(\mathfrak g)\)。
\end{definition}
\symidx{\(Z(\mathfrak g),\operatorname{Der}(\mathfrak g)\)：Lie 代数的中心与导子代数}
中心是一个交换理想。对一个子空间 \(U\subseteq\mathfrak g\)，还分别定义
\[
C_{\mathfrak g}(U)=\bigl\{\,x\mid[x,U]=0\,\bigr\},\qquad
N_{\mathfrak g}(U)=\bigl\{\,x\mid[x,U]\subseteq U\,\bigr\},
\]
称为中心化子与正规化子。前者要求与 \(U\) 中每个元素取括号都为零，后者只要求取括号后仍在 \(U\) 中。
\symidx{\(C_{\mathfrak g}(U),N_{\mathfrak g}(U)\)：Lie 代数中子空间的中心化子与正规化子}
\begin{proposition}\label{b5:prop:lie-centralizer-normalizer}
设 \(\mathfrak g\) 为域 \(k\) 上的 Lie 代数，\(U\subseteq\mathfrak g\) 为任意线性子空间。则中心化子 \(C_{\mathfrak g}(U)\) 与正规化子 \(N_{\mathfrak g}(U)\) 都是 Lie 子代数。若 \(U\) 本身为子代数，则 \(N_{\mathfrak g}(U)\) 是包含 \(U\) 且使 \(U\) 成为理想的最大 Lie 子代数。
\end{proposition}
\begin{proof}
两个条件都对线性组合封闭，所以先得到线性子空间。对 \(x,y\in N_{\mathfrak g}(U)\)、\(u\in U\)，Jacobi 恒等式给出
\[
[[x,y],u]=[x,[y,u]]-[y,[x,u]].
\]
内层的两个括号都在 \(U\) 中，再分别与 \(x,y\) 取括号仍在 \(U\) 中，故 \([x,y]\in N_{\mathfrak g}(U)\)。若 \(x,y\in C_{\mathfrak g}(U)\)，内层两个括号都为零，同一公式便给出 \([x,y]\in C_{\mathfrak g}(U)\)。这些计算都不要求 \(U\) 对括号封闭。

若 \(U\) 为子代数，则 \([U,U]\subseteq U\)，所以 \(U\subseteq N_{\mathfrak g}(U)\)；正规化子的定义又说明 \(U\) 是其中的理想。若另有子代数 \(\mathfrak h\) 包含 \(U\)，并使 \(U\) 成为理想，则 \([\mathfrak h,U]\subseteq U\)，因而 \(\mathfrak h\subseteq N_{\mathfrak g}(U)\)。这就是所述最大性。
\end{proof}

\begin{proposition}\label{b5:prop:lie-derivation-first-order}
设 \(\mathfrak g\) 为域 \(k\) 上的 Lie 代数，\(D\in\operatorname{End}_k(\mathfrak g)\)，令 \(R=k[\varepsilon]/(\varepsilon^2)\)。在 \(\mathfrak g_R=\mathfrak g\otimes_kR\) 上将括号与 \(D\) 分别作 \(R\)-双线性和 \(R\)-线性延拓。则 \(I+\varepsilon D\) 是可逆的 \(R\)-线性映射，且它保持 Lie 括号当且仅当 \(D\) 是 \(\mathfrak g\) 的导子。
\end{proposition}
\begin{proof}
因 \(\varepsilon^2=0\)，\((I+\varepsilon D)(I-\varepsilon D)=I\)，逆映射为 \(I-\varepsilon D\)。对 \(x,y\in\mathfrak g\)，两边分别为
\[
\begin{aligned}
(I+\varepsilon D)[x,y]
&=[x,y]+\varepsilon D([x,y]),\\
[(I+\varepsilon D)x,(I+\varepsilon D)y]
&=[x,y]+\varepsilon\bigl([D(x),y]+[x,D(y)]\bigr).
\end{aligned}
\]
作为 \(k\) 向量空间，\(\mathfrak g_R=\mathfrak g\oplus\varepsilon\mathfrak g\)，所以比较 \(\varepsilon\) 的系数，恰好得到导子等式。保持括号的等式对两个输入均为 \(R\) 线性，故在 \(x,y\in\mathfrak g\) 上成立就能延拓到所有 \(\mathfrak g_R\) 中的元素。
\end{proof}
这说明导子的 Leibniz 条件来自保持括号的自同构的一次项，不必先引入微分或 Lie 群。双数环把二次及更高次项置零，使这一阶条件在任意特征下都是精确等式。
\begin{proposition}\label{b5:prop:lie-derivations}
设 \(\mathfrak g\) 为域 \(k\) 上的 Lie 代数。导子空间 \(\operatorname{Der}(\mathfrak g)\) 在算子交换子下组成 Lie 代数。对每个 \(x\in\mathfrak g\)，映射
\[
\operatorname{ad}x:\mathfrak g\longrightarrow\mathfrak g,\qquad
y\longmapsto[x,y]
\]
为导子，并对每个 \(D\in\operatorname{Der}(\mathfrak g)\) 满足
\[
[D,\operatorname{ad}x]=\operatorname{ad}(D x).
\]
因此这些内导子的空间 \(\operatorname{ad}\mathfrak g\) 为 \(\operatorname{Der}(\mathfrak g)\) 的理想。
\end{proposition}
\symidx{\(\operatorname{ad}x\)：Lie 代数元素的伴随作用}
\begin{proof}
对导子 \(D,E\)，两次使用导子法则得到
\[
\begin{aligned}
DE([x,y])&=[DEx,y]+[Ex,Dy]+[Dx,Ey]+[x,DEy],\\
ED([x,y])&=[EDx,y]+[Dx,Ey]+[Ex,Dy]+[x,EDy].
\end{aligned}
\]
相减后两个混合项分别消去，得到 \([[D,E]x,y]+[x,[D,E]y]\)，所以交换子仍为导子。Jacobi 的导子形式直接说明 \(\operatorname{ad}x\) 为导子。再对 \(y\) 计算
\(D\bigl([x,y]\bigr)-\bigl[x,D(y)\bigr]=[Dx,y]\)，得到最后的等式。
\end{proof}
导子未必都为内导子。例如一维交换 Lie 代数的全部内导子为零，每个线性自同态却都是导子。在任意域 \(k\) 上，取\cref{b5:exm:abelian-heisenberg-lie}中的 Heisenberg 代数 \(\mathfrak h\)，其基为 \(x,y,z\)，非零基括号为 \([x,y]=z\) 及其反向。定义
\[
D(x)=x,\qquad D(y)=0,\qquad D(z)=z.
\]
在 \([x,y]=z\) 上，导子等式为 \(z=[x,y]+[x,0]\)；在 \([x,z]=[y,z]=0\) 上，两边也为零。由双线性与交替性，这已验证全部括号，所以 \(D\) 是导子。但每个 \(u=ax+by+cz\) 都有 \(\operatorname{ad}u(x)=-bz\)，不能等于 \(D(x)=x\)，故 \(D\) 不是内导子。

\ProseParagraphBreak

这个导子也能从一个实际的缩放族看见。对任意 \(t\in k^\times\)，规定
\[
\Phi_t(x)=tx,\qquad \Phi_t(y)=y,\qquad \Phi_t(z)=tz.
\]
则 \([\Phi_t(x),\Phi_t(y)]=tz=\Phi_t(z)\)，其他基括号仍为零，且 \(\Phi_{t^{-1}}\) 是逆映射，故 \(\Phi_t\) 为 Lie 自同构。在双数环中取单位参数 \(t=1+\varepsilon\)，同一公式正好成为 \(\Phi_{1+\varepsilon}=I+\varepsilon D\)。因此这里的一阶缩放与\cref{b5:prop:lie-derivation-first-order}中的导子条件直接相符。

\subsection{伴随表示}
\label{b5:subsec:1102:03}

\begin{proposition}\label{b5:prop:adjoint-representation}
设 \(\mathfrak g\) 为域 \(k\) 上的 Lie 代数。映射
\[
\operatorname{ad}:\mathfrak g\longrightarrow\mathfrak{gl}(\mathfrak g),
\qquad x\longmapsto\operatorname{ad}x
\]
是 Lie 代数同态，其核为 \(Z(\mathfrak g)\)，并诱导 Lie 代数同构
\[
\mathfrak g/Z(\mathfrak g)\cong\operatorname{ad}\mathfrak g.
\]
一个子空间 \(\mathfrak a\subseteq\mathfrak g\) 在所有 \(\operatorname{ad}x\) 下不变，当且仅当它是理想。
\end{proposition}
\begin{proof}
对 \(z\in\mathfrak g\)，Jacobi 恒等式给出
\[
[\operatorname{ad}x,\operatorname{ad}y](z)
=\bigl[x,[y,z]\bigr]-\bigl[y,[x,z]\bigr]
=\bigl[[x,y],z\bigr].
\]
故它保持括号。\(\operatorname{ad}x=0\) 恰好表示 \(x\) 与全部元素取括号都为零，故核为 \(Z(\mathfrak g)\)，像为 \(\operatorname{ad}\mathfrak g\)。对这一同态应用\cref{b5:prop:lie-isomorphism}，所得同构具体为
\[
x+Z(\mathfrak g)\longmapsto\operatorname{ad}x.
\]
最后，\(\mathfrak a\) 在每个 \(\operatorname{ad}x\) 下不变，正是 \([\mathfrak g,\mathfrak a]\subseteq\mathfrak a\)，即理想条件。
\end{proof}
这称为\term[bansui-biaoshi]{伴随表示}[adjoint representation]。它把 Lie 代数的内部括号转成一组线性算子，因而可用特征值、迹与三角化研究结构。取上述任意域上的三维 Heisenberg 代数 \(\mathfrak h\)，对 \(u=ax+by+cz\) 有
\[
[u,x]=-bz,\qquad [u,y]=az,
\]
所以 \(Z(\mathfrak h)=kz\)。同时
\[
\operatorname{ad}x(y)=z,\qquad
\operatorname{ad}y(x)=-z,\qquad
\operatorname{ad}z=0.
\]
前两个算子都非零；若 \(a\operatorname{ad}x+b\operatorname{ad}y=0\)，先作用于 \(y\) 得 \(a=0\)，再作用于 \(x\) 得 \(b=0\)，故伴随像为二维。但
\[
[\operatorname{ad}x,\operatorname{ad}y]
=\operatorname{ad}[x,y]=\operatorname{ad}z=0,
\]
这两个算子彼此交换。因此非交换的 \(\mathfrak h\) 的伴随像却是二维交换 Lie 代数，恰好同构于 \(\mathfrak h/kz\)。一般地，伴随像交换当且仅当 \([\mathfrak g,\mathfrak g]\subseteq Z(\mathfrak g)\)：算子的交换子为 \(\operatorname{ad}[x,y]\)，而 \(\operatorname{ad}\) 的核正是中心。模去中心不仅消去中心方向，也可能使原来的非零括号变为零。

\ProseParagraphBreak

例如二维代数取基 \(h,e\)，规定 \([h,e]=e\)。在这组基下，
\[
\operatorname{ad}h=\begin{pmatrix}0&0\\0&1\end{pmatrix},
\qquad
\operatorname{ad}e=\begin{pmatrix}0&0\\-1&0\end{pmatrix}.
\]
两矩阵的交换子等于第二个矩阵，与原来的关系一致。这个小例子以后将区分可解性与幂零性。
